\section{Reversing the Modern Spell}

Cologero writes:

\begin{quotex}
There is nothing here to oppose effective rites and rituals, or even the acts of a traditional science. After all, the physical world itself is a symbol, a reflection, of the metaphysical. But that is a topic in itself.
\end{quotex}


From the writings of the UR Group:

\begin{quotex}
Yet it would be useless to seek a parallel between the organs and functions of the physical body on the one hand and the inner essence of man on the other hand, since the former are determined by conditions proper to the animal life and by their relationship with the external world. Thus, they represent a deviation, albeit one necessary to certain purposes of existence. Therefore, we cannot directly conclude from the function of an organ, as known to ordinary consciousness, its value as a symbol and expression of the Inner Man.
\flright{\textsc{Leo}, \emph{First Steps Toward the Experience of the Subtle Body}}
\end{quotex}


Much could be commented here, in regards to both of these glosses. We will content ourselves with noting that there is a certain implicit theory of \emph{Auctoritas} inherent in both of these views, which are the same truth expressed from two different facets of that truth. But since most people are ``stuck'' in ordinary consciousness as an existential fact, they would see them as contradictory, and therefore in order to transcend or overcome this, they would have to rely on someone else to point their way towards the truly real. They would have to ``believe'' and then do more than just believe, that both of these statements were true, not because Truth is contradictory, but because they are both obviously woven from the same thread.

Thus, the necessity of Dante, and other ``interpreters'' who can ``read the runes''. Much of what the modern world enacts politically or otherwise is an effective reversal of the proper order of Space and Time\footnote{\url{http://www.gornahoor.net/?p=3166}}. Man was built to re-unify Time, and to divide Space, thus granting himself immortality under God, with ``many mansions'' to inhabit. Modern man concentrates on inverting this pattern -- chopping Time up into meaningless segments of factoids, and throwing everything into a mulligan stew to simmer until the entire planet boils over. In fact, modern man thinks there is no pattern there at all\footnote{\url{https://en.wikipedia.org/wiki/Orders_of_creation}} -- the illusion is reality. Thus, he can do as he pleases, or so imagines.

What we can say about Symbol and Materiality is this -- the link between the two, like the silver cord or golden thread which links the various bodies which are possible of realization in man, is where one can find the ``answers''. It is the place where heaven and earth ``kiss'' (as symbolized in Dante) that holds the interest of one who yearns to follow Romanity, Tradition, \& Reality.

If finding that link is difficult and confusing, this is all the more reason for Gornahoor to emphasize the practical, experiential, existential necessity to do so, provided one builds with sound
\emph{theoria}, rather than experimenting in the modern \emph{ad hoc} fashion, aiming at nothing more than ``something new''.

While it is certainly true that it is difficult to pinpoint true North (witness the many criticisms of Gornahoor made from various angles), this does not mean that it is impossible. After all, God has sent many prophets, messengers, angels, witnesses (among whom was Dante), all of whom will answer a specific yearning of the true needle of the heart, as it homes Northward. With God, and with the heart made in His image, nothing is impossible.


\flrightit{Posted on 2012-03-31 by Logres}

\begin{center}* * *\end{center}

\begin{footnotesize}
\begin{sffamily}
\texttt{Cologero on 2012-04-01 at 13:18 said:}

Where exactly do you see ``many criticisms of Gornahoor made from various angles'', at least of any merit? Rather, there is near total silence, and not the Silence that de Giorgio refers to.

Why is true North so hard to find? Hasn't the compass been invented yet, or do you prefer not to use it? The alternative is to engage in perpetual war with neither end nor resolution, like the Highlanders. We know the prize, but who has the courage to claim it?

\hfill

\texttt{Logres on 2012-04-01 at 20:48 said:}

It's difficult for some, apparently, particularly when it comes to letting go of the material and seeing what is behind it; the only merit I claim is that I recognize this, and repent -- I think again. Something being difficult of application doesn't make it unattainable, if the person struggling with it is able to reorient themselves. Or does it?

\hfill

\texttt{Logres on 2012-04-01 at 21:04 said:}

I've not read any of any merit, as of yet.

\hfill

\texttt{Yakubu Abdul-Aziz on 2012-04-23 at 11:27 said:}

I have the magical spell book by Anna Riva, i dont understand most of the things in it but i want to use it.

\hfill

\texttt{logres on 2012-04-23 at 23:03 said:}

If you are ``magically'' inclined, you should consider Evola's work on the subject:
\url{http://www.amazon.com/Introduction-Magic-Rituals-Practical-Techniques/dp/0892816244}

Theologically, Palamidessi \& Evola are safe guides to such territory. I have no idea what Anna Riva's ideas or praxis would look like. Does it attempt to transcend the spirit world, or merely to dominate it?

\end{sffamily}
\end{footnotesize}

